\SourceLeaf{399}%
\[
\begin{array}{ll}
x+n\tfrac12(x'+x''-2x),&y+n\tfrac12(y'+y''-2y)\\
x'+n\tfrac12(x''+x-2x'),&y'+n\tfrac12(y''+y-2y')\\
x''+n\tfrac12(x+x'-2x''),&y''+n\tfrac12(y+y'-2y'')
\end{array}
\]
oder wenn man
\[
\begin{aligned}1-n&=\alpha\\\tfrac12n&=\beta\end{aligned}
\]
setzt
\[
\begin{array}{ll}
\alpha x+\beta(x'+x''),&\alpha y+\beta(y'+y'')\\
\alpha x'+\beta(x''+x),&\alpha y'+\beta(y''+y)\\
\alpha x''+\beta(x+x'),&\alpha y''+\beta(y+y')
\end{array}
\]
\par
Von den Punkten \(C,C',C''\) werden Perpendikel auf \(A'A'',A''A,AA'\), gefällt, man sucht die Lage der drei Durchschnittspunkte dieser Perpendikel. Es seien die Coordinaten des Durchschnittspunktes der beiden letzten Perpendikel
\[
\xi,\quad\eta
\]
welche man mit Hülfe des folgenden Lehrsatzes bestimmen wird.

Wenn
\[
\begin{array}{ll}
a,&b\\a',&b'\\a'',&b''\\a''',&b'''
\end{array}
\]
die Coordinaten von vier Punkten sind, und die geraden Linien durch den ersten und zweiten Punkt auf der Linie durch den dritten und vierten senkrecht sind, so hat man
\[
\begin{aligned}
\tfrac{b'-b}{a'-a}&=\text{tang. der Neigung der ersten Linie gegen die Abscissenlinie}\\
\tfrac{b'''-b''}{a'''-a''}&=\text{tang. der Neigung der zweiten Linie gegen die Abscissenlinie}
\end{aligned}
\]
und folglich, da die eine Neigung um \(90^0\) grösser ist als die andere, das Product der beiden Tangenten \(=-1\), also
\[
\tfrac{b'-b}{a'-a}.\tfrac{b'''-b''}{a'''-a''}=-1
\]
In unserm Falle hat man also
